% !PRINTED 89
\weberpage{89}
%<<CONT

\begin{equation}
 h=\nu+1,\nu+2,\ldots,n,
 \qquad
 k=\nu+1,\nu+2,\ldots,m,
\tag{4}
\end{equation}
also
\begin{equation}
 (n-\nu)(m-\nu),
\tag{5}
\end{equation}
der Zahl nach.

Diese Bedingungen sind aber wirklich von einander unabhängig, d.~h. es folgt keine aus den übrigen; denn die linken Seiten von (3) können durch geeignete Annahmen über die Coëfficienten $a$ für jede Indexcombination aus der Reihe (4) einen ganz beliebigen Werth erhalten, wie man erkennt, wenn man sämmtliche $a_h^{(i)},a_\mu^{(k)}$ mit Ausnahme von $a_h^{(k)}$ gleich Null setzt.



\begin{center}
{\large \S.\ 26.}

\vspace{0.5em}
{\large\gesperrt{Unhomogene lineare Gleichungen.}}
\end{center}

Wir haben die Aufgabe der Auflösung linearer Gleichungen in den bisherigen Betrachtungen dadurch nicht unwesentlich vereinfacht und auf allgemeinere Gesetze zurückgeführt, dass wir die Gleichungen in Bezug auf die Unbekannten homogen vorausgesetzt haben. In den Anwendungen kommen aber häufig die Unbekannten nicht homogen vor, und wenn auch principiell der eine Fall von dem anderen nicht wesentlich verschieden ist, so wollen wir doch den Fall der nicht homogenen Gleichungen noch besonders betrachten.

Wir können ihn aus dem Fall der homogenen Gleichungen dadurch ableiten, dass wir die Forderung hinzufügen, eine bestimmte der Unbekannten soll den Werth $1$ haben. Wenn diese Unbekannte unter denen vorkommt, die das Problem willkürlich lässt, so entspringt daraus gar keine Schwierigkeit, weil wir sie ja nur $=1$ zu setzen brauchen. Gehört sie aber zu denen, die durch die übrigen bestimmt sind, so müssen noch gewisse Bedingungen erfüllt sein, die besagen, dass der Werth $1$ für diese Unbekannte zulässig ist.

Wir wollen hier die Frage selbständig und in etwas geänderter Bezeichnung behandeln, beschränken uns aber der Einfachheit halber auf den wichtigsten Fall, wo die Anzahl der Unbekannten mit der Anzahl der Gleichungen übereinstimmt. Es sei folgendes System von Gleichungen in Bezug auf die Unbekannten $x_1,x_2,\ldots,x_n$ aufzulösen:
%CONT>>
